%! Discrete Random Variables — Unit 4, Lesson 4.1 — Slides (Beamer)
\documentclass[aspectratio=169, 11pt]{beamer}
\usepackage{saar-beamer}   % provides \ceruleanheader{} and \sectionlabel[color]{}

\begin{document}

% TITLE SLIDE (CourseName is not defined in beamer, so the name is literal)
{
\setbeamercolor{background canvas}{bg=cerulean}
\begin{frame}[plain]
  \vspace{2.8cm}\hspace{0.55cm}%
  \begin{minipage}{0.9\textwidth}
    {\color{goldacc}\bfseries\small LESSON 4.1}\\[0.5cm]
    {\color{white}\bfseries\LARGE Discrete Random Variables}\\[1.2cm]
    {\color{frostmid}\small Statistical Analysis \& Algebraic Reasoning~$\cdot$~Unit 4}
  \end{minipage}
\end{frame}
}

% The deck follows the EFFL flow (~9 frames):
%   learning targets (SPOILER-FREE) → warm-up → activity launch → 3-4 debrief frames
%   (formalize in "red ink") → QuickNotes summary → application → close & assign.
% Everything before the debrief frames stays vocabulary-free (the spoiler rule).
\newcommand{\redink}[1]{{\color{redacc}\bfseries #1}}

% LEARNING TARGETS (spoiler-free)
\begin{frame}[t, plain]
  \ceruleanheader{Today: experience first, formalize later}
  \sectionlabel{Learning targets}
  By the end of today, I can\ldots
  \begin{itemize}\setlength{\itemsep}{6pt}
    \item TODO: plain-language target (no formal vocabulary).
  \end{itemize}
  \begin{block}{How today works}
    Warm-up $\to$ \textbf{TODO activity} with your group $\to$ we name what you found $\to$
    one application together $\to$ practice.
  \end{block}
\end{frame}

% WARM-UP
\begin{frame}[t, plain]
  \ceruleanheader{Warm-Up: things you already know}
  \sectionlabel{5 minutes --- alone, then check with a neighbor}
  % TODO: the warm-up items.
\end{frame}

% ACTIVITY LAUNCH
\begin{frame}[t, plain]
  \ceruleanheader{Experience \& Formalize: TODO Activity Title}
  \sectionlabel[goldacc]{Groups of 3--4 $\cdot$ 22 minutes}
  % TODO: the data context + scenarios, and a block:
  \begin{block}{Your job}
    Use what you already know. For \emph{every} answer, be ready to show me
    \textbf{where you see it in the data}.
  \end{block}
\end{frame}

% DEBRIEF (repeat this frame 3-4 times; formal terms in \redink{})
\begin{frame}[t, plain]
  \ceruleanheader{Debrief: naming what you found}
  \sectionlabel[redacc]{TODO}
  % TODO: ``student phrase'' $\to$ \redink{formal term} moves.
\end{frame}

% QUICKNOTES SUMMARY
\begin{frame}[t, plain]
  \ceruleanheader{QuickNotes: TODO topic}
  \sectionlabel{Fill the box in your packet}
  % TODO: the worked example the QuickNotes box records.
\end{frame}

% APPLICATION
\begin{frame}[t, plain]
  \ceruleanheader{Application: TODO title}
  \sectionlabel[goldacc]{We work this one together --- you hold the pen}
  % TODO: the problem, then the interpret-in-context question and the "what if" that
  % demands a justification. This is the lesson's only in-class practice.
\end{frame}

% CLOSE & ASSIGN
\begin{frame}[t, plain]
  \ceruleanheader{Practice \& what's next}
  \sectionlabel[goldacc]{Check Your Understanding --- this one is scored}
  % TODO: name the assignment and say plainly which form it takes today — the packet's
  % Check Your Understanding pages OR the equivalent DeltaMath set. Give the due date and
  % how it is scored; close with a one-line preview of the next lesson.
\end{frame}

\end{document}
